\begin{figure}[H]
\centering
\includegraphics[width=1.0\linewidth]{ehime/M5_drift_month_grid.pdf}
\caption[月別の平均落下位置]{各格子点から上げた場合の平均落下位置。\textbf{矢印は実寸}で、矢先がそのまま予測落下点である。色が平均の飛行距離。全パネルが同じ枠と同じ色スケールなので、月どうしを直接比べられる。}
\label{fig:m5-ehime}
\end{figure}

\begin{figure}[H]
\centering
\includegraphics[width=0.92\linewidth]{ehime/M7-01_drift_map.pdf}
\caption[1 月の落下位置]{1 月前半の平均落下位置を拡大したもの。枠は集計対象の領域で、矢先はその外へ出ていることが多い。}
\label{fig:m701-ehime}
\end{figure}

\begin{figure}[H]
\centering
\includegraphics[width=0.92\linewidth]{ehime/M7-07_drift_map.pdf}
\caption[7 月の落下位置]{同じく 7 月前半。1 月と比べると距離も向きも大きく違う。}
\label{fig:m707-ehime}
\end{figure}

\begin{figure}[H]
\centering
\includegraphics[width=1.0\linewidth]{ehime/M6_drift_percentile_map.pdf}
\caption[荒れた日の飛行距離]{飛行距離の 90 パーセンタイル。平均 \SI{150}{\kilo\metre} で p90 が \SI{400}{\kilo\metre} という状況と、平均 \SI{150}{\kilo\metre} だけを見て計画することは、運用上まったく違う。}
\label{fig:m6-ehime}
\end{figure}

\begin{figure}[H]
\centering
\includegraphics[width=0.92\linewidth]{ehime/M8_landing_ellipse_map.pdf}
\caption[代表地点からの落下分布]{1 地点から上げた場合の落下位置を季節別に重ねたもの。楕円は 50 / 90 / 95 \% の範囲で、各解析時刻ごとに求めた着地点の共分散から描いている。散らばりが 2 変量正規に従うという仮定は残る。灰色が陸、白が海である。}
\label{fig:m8-ehime}
\end{figure}
